\begin{tikzpicture}[x=1cm,y=1cm]
  % ---------- step 1: the 2 x 4 tableau of one Apply microstep at F=1
  \node[lbl,anchor=east] at (-0.15,-0.35) {\(C_0\)};
  \foreach \j/\txt in {0/{$\langle q_0,\mathsf{Ap}\rangle$}, 1/{$\lambda$},
                       2/{$v$}, 3/{$\sqcup$}}{
    \node[cell, minimum width=1.55cm, minimum height=7mm, font=\scriptsize,
          anchor=north west] (a\j) at (\j*1.55,0) {\txt};
  }
  \node[lbl,anchor=east] at (-0.15,-1.55) {\(C_1\)};
  \foreach \j/\txt in {0/{$\mathsf{ret}$}, 1/{$\langle q_1,v\rangle$},
                       2/{$v$}, 3/{$\sqcup$}}{
    \node[cell, minimum width=1.55cm, minimum height=7mm, font=\scriptsize,
          anchor=north west] (b\j) at (\j*1.55,-1.2) {\txt};
  }
  \node[draw=focus, line width=1.1pt, rounded corners=1.5pt,
        fit=(a0)(a2), inner sep=1.2pt] (win) {};
  \node[cell head, minimum width=1.55cm, minimum height=7mm, font=\scriptsize,
        anchor=north west] at (1.55,-1.2) {$\langle q_1,v\rangle$};
  \node[lblsm,text=term,anchor=south west] at (0,0.12)
    {radius-one window \((j{-}1,j,j{+}1)\) at \(j=2\)};
  \draw[->,machine,line width=1pt] (2.33,-0.78) -- (2.33,-1.2)
    node[midway,right=1mm,lblsm,text=machine] {\(\omega\)};
  \node[lbl,anchor=west,align=left,text width=6.6cm] at (6.5,-0.9)
    {\textbf{one \(\operatorname{Apply}\) microstep.}\ \ The term
     \(\operatorname{Apply}(\operatorname{Lambda}(1,\operatorname{BoundVar}(0,0)),v)\)
     at fuel \(F=1\) has a \(2\times4\) tableau: two rows, width four.  The
     amber window of row \(C_0\) fixes the amber cell of row \(C_1\).};

  % ---------- step 2: one clause index of the succinct 3SAT circuit
  \node[descbox, anchor=north west, minimum width=4.6cm, font=\scriptsize\sffamily]
       (idx) at (0,-2.55) {index \((\mathrm{transition},\,i=0,\,j=2)\)};
  \node[checked, anchor=north west, minimum width=7.4cm, font=\small]
       (cl3) at (5.4,-2.55)
    {\(\neg y_w\ \vee\ x_{1,2,s}\ \vee\ x_{1,2,s}\)};
  \draw[xform] (idx.east) -- (cl3.west);
  \node[lblsm,anchor=north west,align=left,text width=12.8cm] at (0,-3.55)
    {Cells are numbered from \(1\), so the window at \(j=2\) is cells
     \(1,2,3\): \(y_w\) is the Tseitin variable for ``row \(0\) shows
     \((\langle q_0,\mathsf{Ap}\rangle,\lambda,v)\) at this window'', and
     \(x_{1,2,s}\) says ``row \(1\),
     cell \(2\) carries \(s=\langle q_1,v\rangle\)''.  A two-literal clause is
     padded to three by repetition, as Theorem~\ref{thm:l-succinct-sat}
     requires.};

  % ---------- step 3: the same clause after decoupling
  \node[checked, anchor=north west, text width=12.2cm, align=center, font=\small]
       (cl5) at (0,-5.00)
    {\(a_{i_1}^{o_1}\ \vee\ b_{i_2}^{o_2}\ \vee\ (u_3)^{0}_{\,y_w}\ \vee\
       (u_4)^{1}_{\,x_{1,2,s}}\ \vee\ (u_5)^{1}_{\,x_{1,2,s}}\)\\[1pt]
      {\scriptsize\sffamily for every \(i_1,i_2\) and every
       \(o_1,o_2\in\{0,1\}\)}};
  \draw[xform] (6.4,-4.55) -- (6.4,-5.00);
  \node[tag check, anchor=north east] at (12.7,-6.40) {ONE BIT PER BLOCK};
  \node[lblsm,anchor=north west,align=left,text width=12.8cm] at (0,-6.70)
    {The first, second and third literal go to \(u_3,u_4,u_5\).  Because both
     polarities \(o_1,o_2\) occur for every pair \(i_1,i_2\), some clause of the
     family has both answer literals false, so the family is equivalent to the
     original three-literal clause --- which is why decoupling costs nothing but
     the equality gadgets \(u_3=u_4=u_5\).};
\end{tikzpicture}
